\section{Lilac}
\label{lilac}
The following are a few core intuitions behind Lilac. Lilac aims to reason about random variables, which is at the level of abstraction natural for
human reasoning (as opposed to descending down into measure theory and lebesgue integrals). Lilac's notion of separation is completely different to that
of classical separation logic. The resource here, instead of heaps, are Probability spaces:

\begin{leancode}
structure PSpace (Ω : Type*) where
  ms : MeasurableSpace Ω
  μ : Measure[ms] Ω
  is_prob : IsProbabilityMeasure μ
\end{leancode}

Essentially this just a Probability Measure, a measure which assigns measure (or probability) 1 to the universal set, and
all the data required to define it (a \texttt{MeasurableSpace}). So by analogy with classical separation logic, a Probability
Space dictates the state of a probabilistic program. This is achieved by the values of variables in a probabilistic
program being modelled by based on this probability measure. Concretely, variables in the logic are measurable functions (
aka random variables) from the \( \Omega \marrow ty.den \) for some \( ty : Ty \). Which means that the distribution (the measure) over \( ty.den\) is a 
pushforward of \(\mu\) from the resource. Hence, the resource is controlling the distribution of every random variable which makes up expressions which
makes up probabilistic programs.

\subsection{Separation is probabilistic independence}
Now we can present the notion of separation defined over probability spaces. One of the slogans is ``separation is probabilistic independence''.
probability spaces admit a binary operation, say \( p \star q \) which can guarantee that the random variables defined through the resource p and those constructed through
q are independent. The definition of this binary operation is in a sense a generalisation of independence of random variables: (note we use $\mathcal{E}$ $\mathcal{F}$ and $\mathcal{G}$ for
measurable spaces)

\textbf{Definition 2.2 \citep{lilac} (Independent Combination).}
Let $(\Omega, \mathcal{E},\mu)$ and $(\Omega, \mathcal{F},\nu)$ be probability spaces over the same ambient sample space $\Omega$. 
A probability space $(\Omega,\mathcal{G},\rho)$ is an independent combination of $(\Omega,\mathcal{E},\mu)$ and $(\Omega,\mathcal{F},\nu)$ if (1) $\mathcal{G}$ is 
the smallest $\sigma$-algebra containing $\mathcal{E}$ and $\mathcal{F}$, and (2) $\rho$ witnesses the independence of $\mu$ and $\nu$ in the sense that 
for all $E \in \mathcal{E}$ and $F \in \mathcal{F}$ it holds that $\rho(E \cap F) = \mu(E)\nu(F)$.

The 2nd condition is familiar since it's the standard definition of independence of two events E and F. The novelty is that, independence is lifted from single events to all possible events
occurring in two different measurable spaces. And that's how we get separation; we can just say an arbitrary random variable from p and q respectively are independent.
Remember that we desire a partial binary operation to define a kripke resource monoid (KRM). It is non-trivial to prove that the independent combination
when exists is unique, required to show that independent combination is a function not a relation. The ``pi-lambda-theorem'' to prove both this and several other laws of KRM are satisfied.

\subsection{Measurable Spaces and Random Variables on the Hilbert Cube}
Probability spaces are defined parametrically over some type $\Omega$, but Lilac makes a specific choice for $\Omega$ in order to define
the notion of a finite footprint. We will return to finite footprint later in \ref{}, but here we first define $\Omega = [0, 1]\^{\mathbb{N}}$.
Henceforth denote the interval, $[0,1]$ as $I$.
This type is known as the Hilbert Cube, because the type $\mathbb{R}^{\mathbb{N}}$ is more well known Hilbert Space. The Hilbert Cube is an countably infinite sequence of intervals $I$,
in other words $\mathbb{N} \to I$. First we must establish a measurable space over $I\^{\mathbb{N}}$, to turn this into a probability space.
Intuitively this would be obtained by taking the standard borel $\sigma$-algebra, $\Sigma_I$ over each dimension and taking the product measurable space iteratively over all dimensions.
In other words a fold over $\mathbb{N}$. Mathlib offers a more general way to construct this measurable space, as it defines \texttt{Measurable.pi}, which can construct
does not require a countable index, but works on any $\Pi$-type (function type). Like so:

\begin{leancode}
instance MeasurableSpace.pi {X : δ → Type*} [m : ∀ a, MeasurableSpace (X a)] : MeasurableSpace (∀ a, X a) := ...
\end{leancode}

This is the standard borel $\sigma$-algebra, $\Sigma_{I^\mathbb{N}}$ on the Hilbert Cube. Each \emph{measurable rectangle} is a cartesian product over infinite dimensions. And the general measurable set is
constructed using unions or complements of the rectangles. This is the finest (or largest) Measurable Space over the Hilbert cube that we will work with. Remember we have one measurable
space associated to each probability space, so there will be many different measurable spaces constructed over the Hilbert cube, but they will all be coarser (smaller) than the
$\Sigma_{I^\mathbb{N}}$.
Motivated by the aim of reasoning about random variables, $ \Omega \marrow A$ (where $A$ is some \texttt{ty.den}) we define:
$RV A \triangleq I^{\mathcal{N}} \marrow A \text{such that the map is finite footprint}$. At the moment the finite footprint is an additional restriction that does not
concern us. Measurable Spaces are implicit in the standard mathematical notation, so note that RV is defined w.r.t $\Sigma_{I^\mathbb{N}}$ (\texttt{MeasurableSpace.pi}) on $I^{\mathbb{N}}$.



Remember that $\mathcal{G} A$ is the space of measures over A.

\subsection{Ownership is measurability}
Let's provide some more intuition on how probability spaces can only admit certain random variables (thereby achieving its mutual independence requirements with variables from the other spaces),
This pertains to the 2nd slogan ``ownership is measurability''. To do so, let's introduce the 
syntax and semantics of the logic:

\[
\begin{aligned}
S,T &\in \mathbf{Set} \\
A,B &\in \mathbf{Meas} \\
P,Q ::=~& \top \mid \bot \mid P \land Q \mid P \lor Q \mid P \to Q \mid P \ast Q \mid P \wand Q \mid {} \\
& \forall x\!:\!S.\,P \mid \exists x\!:\!S.\,P \mid \forall_{\mathrm{rv}} X\!:\!A.\,P \mid \exists_{\mathrm{rv}} X\!:\!A.\,P \mid {} \\
& E \sim \mu \mid \mathrm{own}\,E \mid E \circeq E \mid \mathbb{E}[E] = e \mid \mathrm{wp}(M,X\!:\!A.\,Q)
\end{aligned}
\]

Notice that we have two sets of quantifiers: one for deterministic variables and one for random variables.
Over the course of the mechanisation, I noticed that in Core Lilac (without the conditioning modality) does not use \emph{deterministic variables}.
This simplifies our presentation here, omitting deterministic contexts and variables, and the two quantifiers $\forall$ and $\exists$ above. The following
rules are helpful to understand, how to construct terms in the syntax. We have already omitted deterministic context and quantifiers here:

\begin{figure}[hbtp]
\centering
\begin{gather*}
\Delta ::= {} \cdot \mid \Delta, X : A \\[3pt]
\sem{X_1 : A_1, \ldots, X_n : A_n} = A_1 \otimes \cdots \otimes A_n
\end{gather*}
\begin{mathpar}
\inferrule{\Delta \vdash P \\ \Delta \vdash Q}{\Delta \vdash P \oplus Q} \and
\inferrule{\Delta, X : A \vdash P}{\Delta \vdash \forall_{\rv} X{:}A.\, P} \and
\inferrule{\Delta, X : A \vdash P}{\Delta \vdash \exists_{\rv} X{:}A.\, P} \and
\inferrule{E \in \sem{\Delta} \marrow A \\ \mu \in \mathcal{G}A}{\Delta \vdash E \sim \mu} \and
\inferrule{E \in \sem{\Delta} \marrow A}{\Delta \vdash \mathsf{own}\, E} \and
\inferrule{E_1 \in \sem{\Delta} \marrow A \\ E_2 \in \sem{\Delta} \marrow A}{\Delta \vdash E_1 \circeq E_2} \and
\inferrule{E \in \sem{\Delta} \marrow \mathbb{R} \\ e \in \mathbb{R}}{\Delta \vdash \mathbb{E}[E] = e} \and
\inferrule{M \in \sem{\Delta} \marrow \mathcal{G}A \\ \Delta, X : A \vdash Q}{\Delta \vdash \mathrm{wp}(M, X{:}A.\, Q)}
\end{mathpar}
\caption{Simplified Typing rules of Lilac}
\label{fig:typing}
\end{figure}

where we use this notation $\oplus \in \{\land, \lor, \to, *, \wand\}$
% \[
% \begin{array}{@{}c@{\qquad\qquad}c@{\qquad\qquad}c@{\qquad\qquad}c@{\qquad\qquad}c@{}}
% &
% E \in \RV A
% &
% e \in A
% &
% \mu \in \mathcal{G}A
% &
% M \in \RV \mathcal{G}A
% \end{array}
% \]

Lilac like any other separation logic has additional connectives beyond the standard connectives specified by BI that is required of it. The first row of the syntax above standard connectives in BI,
so their semantics is identical that of classical separation logic, replacing heaps with probability spaces. The second row of quantifiers is explained later when explaining iris proof mode, which
requires a different treatment of quantifiers to keep track of variables when carrying out proofs in the logic. The connectives in the third row
are called logic (or probability) specific connectives and where we make interesting assertions about the state of a program. To understand the relation between the logic-specific connectives and the standard connectives, think of the syntax of the logic as a tree,
the logic specific connectives tend to be the leaves, while the standard connectives are like nodes (they are mostly binary). The nodes are there to compose together information expressed by the leaves.

The central probability specific connective is \texttt{Own E}. Fig \ref{fig:prob-connectives} shows the semantics (or \emph{satisfiability relation} in this context).
Notice that both the \(~\) (distributed as) and \( \mathbb{E} \) connectives have the semantics of \texttt{Own E}, \( E \text{ is } \mathcal{F}\text{-measurable} \), with an additional conjunct.
This leads to the intuition that a lilac proposition (henceforth \texttt{LProp}) must ``own'' a random variable before saying anything about it. In probabilistic separation logic we interpret ownership
of a random variable as enforcing that no other separating conjunct may own a random variable which is not independent. But how do we enforce this using the semantics? Why is ownership \texttt{Own E},
measurability \( E \text{ is } \mathcal{F}\text{-measurable} \)? First let us clarify that we already know that $E$ is a measurable function $RV E \triangle I^{\mathbb{N}} \marrow A$ so is
\texttt{Own E} redundant? No. $RV E$ defines measurability w.r.t to $\Sigma_{I^\mathbb{N}},$ where  $\Sigma_{I^\mathbb{N}} \geq \mathcal{F}$, This measurability w.r.t $\mathcal{F}$ is a stronger condition.
To see this simply take some function and notice that it's preimage on a set in the codomain, if contained in $\mathcal{F}$, must also be in $\Sigma_{I^\mathbb{N}}$.

To provide an intuition for why ownership is defined as \( E \text{ is } \mathcal{F}\text{-measurable} \), let's consider a scenario where some dimension in the Hilbert cube, n,
is used to define a random variable. Take $\omega \in I^\mathbb{N} \triangleq \mathbb{N} \to I$. and $X : RV \mathbb{R} := \lambda \omega \mapsto \omega n$. This projection is a random
variable. Now we can study under what measurable spaces of $I^\mathbb{N}$ ``own X'' holds. We split the hilbert cube at dimension n into 3 chunks: $0 \ldots n-1, n, n+1 \ldots \infty$,
and talk about the measurable spaces in each piece. The product measurable spaces of these 3 pieces will form the whole hilbert cube measurable space again.
We do not care at all about what $\sigma$-algebra is present in the first and the last chunk, but only the $\sigma$-algebra at dimension n matters. If at n we have the trivial $\sigma$-algebra,
$\{\phi, I\}$, then X is certainly not measurable. In fact, we cannot have any $\sigma$-algebra smaller than the borel $\sigma$-algebra on $I$, $\Sigma_I$, to validate \texttt{Own X}.

And if two probability spaces $(I^\mathbb{N}, \mathcal{E}, \mu)$,  $(I^\mathbb{N}, \mathcal{F}, \nu)$ have the same $\sigma$-algebra at dimension n, $\mathcal{E}\!\mid_n = \mathcal{F}\!\mid_n$,
Then the independent product of these two spaces will not exist in general. To see this take for a moment the simpler probability spaces $(I, \Sigma_I, \mu\!\mid_n)$ and $(I, \Sigma_I, \nu\!\mid_n)$,
by focussing our attention on the nth dimension. The independent combination would be $(I, \Sigma_I, \lambda x \mapsto \mu\!\mid_n x * \nu\!\mid_n )$. But the last element of that tuple is not
a measure! It would not satisfy the additivity condition on measure that taking the measure of a union of measurable sets is the sum of measure of the sets.
The takeaway is that identical $\sigma$-algebras cannot form independent combinations, and in this example we illustrated how measurability was unique ownership of a dimension of the Hilbert cube.
Measurability over coarser $\sigma$-algebra than $\Sigma_I$ is useful as well: take $Y : RV Bool := \lambda\ \omega \mapsto \omega\ n \le 0.5$. To validate \texttt{Own Y}, a much coarser $\sigma$-algebra on dimension n
than $\Sigma_I$ would suffice: ${\phi, [0,0.5], [0.5,1], I}$. In general measurability can express complex ownership patterns single dimensions can be split across multiple separating conjuncts, and multiple dimensions can be mixed.

\begin{figure}[htbp]
  \centering
  \begin{minipage}{0.9\linewidth}
\[
\begin{aligned}
(\mathcal{F},\mu) \vDash \mathrm{own}\,E &\iff E \text{ is } \mathcal{F}\text{-measurable} \\[0.4em]
(\mathcal{F},\mu) \vDash E \sim \mu' &\iff E \text{ is } \mathcal{F}\text{-measurable} \;\land\; \mu' = E_{\#} \mu \\[0.4em]
(\mathcal{F},\mu) \vDash \mathbb{E}[E] = e &\iff E \text{ is } \mathcal{F}\text{-measurable} \;\land\; \mathbb{E}_{\omega\sim\mu}[E(\omega)] = e \\[0.4em]
(\mathcal{F},\mu) \vDash E_1 \overset{\mathrm{as}}{\equiv} E_2 &\iff F \in \mathcal{F} \;\land\; \mu(F)=1 \;\land\; F \cup (E_1,E_2)^{-1}(A) \in \mathcal{F} \text{ for all } A \in \mathrm{cod}(E_1)\otimes\mathrm{cod}(E_2), \\
&\quad \text{where } F = \{\omega \mid E_1(\omega)=E_2(\omega)\} \\[0.4em]
(\mathcal{F},\mu) \vDash \mathrm{wp}(M,X\!:\!A.\,Q) &\iff \text{for all } \mathcal{P}_{\mathrm{frame}} \text{ and } \mu \text{ with } \mathcal{P}_{\mathrm{frame}}\bullet\mathcal{P} \sqsubseteq (\Sigma_\Omega,\mu) \text{ and all } D_{\mathrm{ext}} : \mathrm{RV}\,\llbracket \Delta_{\mathrm{ext}} \rrbracket, \\
&\quad \text{there exists } X : \mathrm{RV}\,A \text{ and } \mathcal{P}' \text{ and } \mu' \text{ with } \mathcal{P}_{\mathrm{frame}}\bullet\mathcal{P}' \sqsubseteq (\Sigma_\Omega,\mu') \text{ such that } \\
&\quad \left( \begin{array}{l} \omega \leftarrow \mu;\\ v \leftarrow M(\omega);\\ \mathrm{ret}\bigl(D_{\mathrm{ext}}(\omega),v\bigr) \end{array} \right) = \left( \begin{array}{l} \omega \leftarrow \mu';\\ \mathrm{ret}\bigl(D_{\mathrm{ext}}(\omega),X(\omega)\bigr) \end{array} \right), \\
&\quad \text{and } (X),\mathcal{P}' \vDash Q
\end{aligned}
\]
  \end{minipage}
  \caption{Probability Specific Connectives in Lilac}
  \label{fig:prob-connectives}
\end{figure}

\subsection{Variables}
In order to reason about programs, the logic must be able to express facts about the program effectively. In other words, the language and the logic need to 
``communicate'' somehow. In this section we explain the interface between the language and the logic. Interfaces between different components, namely:
language \& logic and logic \& iris, are important in a formalisation as it guides and motivates choices in how key definitions are made.

Generally speaking a logic expresses propositions about a program in the language by, having access to the program variables through some mechanism. Before we go
into this mechanism, there is a novelty in the way Lilac has a different perspective on variables in the language vs the logic. A variable in the language may have
the type \texttt{x : ty} (Appl type judgement) for some \texttt{ty : Ty} (meta language type judgement). One would then expect that if the logic were to construct an
\texttt{LProp} in which this variable occurs, it would have type \texttt{x : ty.den}. However, as we saw above the type of variables in the logic is instead \texttt{X : RV ty.den}.
using this variable. The conversion from \texttt{ty.den} to \texttt{RV ty.den} is a reflection of the perspective taken in probabilistic reasoning. Let's
explain with an example. Take

\[
x \leftarrow \mathrm{unif}\ [0,1];\ y \leftarrow \mathrm{unif}\ [0,1];\ \mathrm{ret}\ (x,y)
\qquad \textsc{(UNIF2)}
\]

\[
\{ \top \}\ \textsc{unif2}\ \{ (X,Y).\ X \sim \mathrm{Unif}[0,1] \ast Y \sim \mathrm{Unif}[0,1] \}
\]

In the program (according to the semantics of bind) $x$ and $y$ are sampled, so they are fixed values. One can consider multiple runs through the program where
the values of the samples are different, but locally $x$ and $y$ are deterministic values. This is in contrast with the logic where $X$ and $Y$ are random variables.
We consider the full distribution of values $x$ and $y$ could have taken across all \emph{runs} of \texttt{UNIF2}.

The link between the language and the logic is established using a connective in the logic, which takes the denotation of a program and some
other assertions regarding it. Lilac uses the \emph{weakest precondition} (wp) connective, but let's first have a look at the more intuitive \texttt{Hoare triple}
connective. It takes 3 arguments, represented syntactically as $\{ P \} \llbracket M \rrbracket \{X. Q \}$.
$P$ (precondition) is an \texttt{LProp} and $\llbracket M \rrbracket$ is the denotation of some program $M$. Q is an \texttt{LProp}
with the (random) variable X free in it. If $M : Ty.G ty$ (Appl type judgement), then $X : RV ty.den$. Since Appl is a functional language, the only variable that
we have access to in the logic is that representing the outcome of the entire program. Though in practice multiple variables can be accessed
by the logic by destructuring a tuple output by the program as in the example of UNIF2.

\Todo{how this destructuring works even with RV?}

The alternative connective, $wp$, is contructed with the denotation of a program and only one LProp, the postcondition. Syntactically: $wp \llbracket M \rrbracket \{X. Q\}$.
$wp$ is defined to be the weakest precondition on the program M, such that the postcondition Q holds. In a consistently defined semantics the following equivalence
constructing a Hoare triple in terms of wp holds: $\{ P \} \llbracket M \rrbracket \{X. Q \} \triangleq P \wand wp \llbracket M \rrbracket \{X. Q \}$. The advantage
of using $wp$ is that it is more amenable to automation and since it only takes in a single proposition. It is easier to infer postconditions in proof rules
which involve going from one wp LProp to another, and makes the prospect of applying these rules easier as less (often no) parameters will need to be
explicitly specified. An example of a rule converting between wp LProps is \texttt{wp\_bind}:

\[
\mathrm{wp}(\llbracket M \rrbracket, X.\,\mathrm{wp}(\llbracket N \rrbracket, Y.\,Q))
\;\vdash\;
\mathrm{wp}(\llbracket X \leftarrow M;\, N \rrbracket, Y.\,Q)
\]

% This causes technical challenges in defining and connecting the language
% and logic in lean and particularly in the instantiating in Iris, since there is no previous instantiation (as far as I am aware) of the variables not being the same across language and logic.

% Note that ``global'' reasoning tends to conflate the variables  does not mean non-compositional still enjoy a form of compositional reasoning using
% the logic, which is the independence of random variables. Thus, the compositionality offered by the semantics of the language, the line-by-line compositionality since
% each term/line is given a denotation which composes cleanly, may be understood as horizontal compositionality. The logic on the other hand, allows us to simultaneously
% talk about the variables introduced in every line of the program through the sequencing (monadic  bind) of programs. However, the structure of assertions implicitly
% inform us of the independence between the variables where present, which can be understood as a \emph{vertical} compositionality.

% \Todo{revise the above into talking about this compositionality stuff elsewhere}

% \Todo{This deep introspection about compositionality goes at the end}
% Since it's a separation logic, it is also compositional but
% along a different axis to compositionality of the denotational semantics. I call the compositionality of probabilistic separation logic a line-by-line compositionality
% (as opposed to term-wise compositionality). Each let binding creates a new "line" and we can compose our reasoning about each of line. In the denotational semantics
% we would instead reason compositionally about the arguments to the monadic bind operator (an s-finite measure and an s-finite-kernel). The former is a more natural way to reason
% about a sequencing operator.

% This reasoning style is achieved by defining a resource which gets carried through and composed and split according to the laws of separation logic. This resource or state is
% a canonical source of randomness, the Hilbert cube, along with a measure space structure over it. It can be argued that the Hilbert cube could be replaced with something else,
% and indeed the definitions are parametric over this choice of measure space for a while. Any distribution can be constructed from the Hilbert cube (which is like a uniform distribution)
% using the operations/terms of the language.




% The hhilbert cube is just a temporary hack ofr defining finite footrpint right??? Wouldn't it be nice to have a principled
% general definition of Finite footprint. Perhaps category theory/ synthetic measure theory can help
